\chapter{Conclusiones y recomendaciones}
\label{cap:conclusiones}
\glsresetall

\section{Conclusiones}

\begin{comment}
    Objetivo Específico 1
    Analizar las tecnologías de cifrado y protocolos de túneles seguros disponibles en el estado del arte, seleccionando la solución que ofrezca el mejor equilibrio entre robustez criptográfica y baja latencia para el protocolo \gls{MODBUSTCP} en el entorno de la Micro-Red.
\end{comment}

El análisis de tecnologías de protección permitió concluir que la solución más adecuada para el enlace evaluado no debía seleccionarse únicamente por su fortaleza criptográfica, sino por su compatibilidad con la operación industrial existente. En el contexto de la Micro-Red, donde la comunicación \gls{SCADA}-\gls{API} utiliza \gls{MODBUS}/\gls{TCP} y donde el controlador no debe asumir gestión de llaves, certificados ni cambios de aplicación, la alternativa seleccionada debía proteger el tráfico sin modificar la semántica del protocolo industrial.


\section{Recomendaciones}

Se recomienda tratar el canal seguro \gls{SCADA}-\gls{API} como un componente permanente de la arquitectura de ciberseguridad de la Micro-Red y no como una configuración aislada del experimento. Para ello, se debe documentar la configuración de interfaces, direcciones, parámetros de \gls{MACsec}, criterios de arranque, parada, diagnóstico y recuperación. También conviene mantener una línea base histórica del tráfico \gls{MODBUS}/\gls{TCP}, de las latencias y del consumo de recursos, de modo que futuras variaciones puedan compararse con evidencia previa y no únicamente con observaciones puntuales.
